% !TeX program = xelatex
\documentclass{UAmathtalk}
\usepackage{setspace}
\renewcommand{\when}{12:00 p.m.}
\renewcommand{\where}{Hudson 0110}
%\renewcommand{\where}{Hudson 0110.\\ \href{https://albany.zoom.us/j/81219478289?pwd=N5fy7vxKBXqlQ4mAq4KaeaAM5iAHi4.1}{Zoom ID: 812\,1947\,8289 - Passcode: 233444}}
\author{ Ken McCabe }
\address{Northeastern University}
\urladdr{https://kenmccabe.net}
\title{Wasserstein Stability of the Nerve Map}
\date{Friday, October 23rd, 2026}


\begin{document}

\maketitle
%\renewcommand*{\abstractsize}{\normalsize}
\renewcommand*{\abstractsize}{\fontsize{13.5}{16.2}\selectfont}
\begin{abstract}
    {
Given a locally finite open cover of a paracompact space, a subordinate partition of unity defines a map into the geometric realization of the nerve, and when the cover is good, the nerve theorem says this map is a homotopy equivalence. In applied topology, one specializes to covers of metric data by balls centered at finitely many landmarks. The corresponding nerve map represents the data inside the \v{C}ech complex, enabling one to pull cohomological or spectral invariants back to the data. In this talk, I will discuss ongoing work with José Perea in which we introduce a weighted generalization of the nerve map and study its stability under perturbations of the landmarks and weights, using metrics from optimal transport. Under a uniform nondegeneracy assumption, we prove that this map is Lipschitz stable with respect to a scale-dependent refinement of the 1-Wasserstein distance, and we show our estimate is sharp. As this distance is dominated by a multiple of the 1-Wasserstein distance itself, we obtain Wasserstein stability as a corollary. We also give explicit estimates of the Lipschitz constants for closed Riemannian manifolds. Finally, time permitting, I will sketch several applications, including to the stability of Laplacian eigenvector cascading and Eilenberg--MacLane coordinates.
    }
\end{abstract}
\end{document}


